\section{Model and two transposition problems}
\label{sec:model}
We use the deterministic multitape model of \cite{FirstPaper2026}, imported
through the shared Lean package of \cite{SecondPaper2026}. A machine has fixed
finite control, a fixed finite work alphabet, and a fixed number of two-way
work tapes. Initially every work cell is blank and every head is at the
origin. Its read-only input is a raw finite bit string; the input end symbol
is visible outside the string. Output is append-only. One transition queries
the scanned symbols, updates the finite control and work symbols, moves each
head by at most one position, and appends at most one output bit. No matrix
header or dimension is supplied separately.

For $d\geq1$, the canonical entry width and input length are
\[
 b(d)=\max\{1,\ceil{\log_2 d}\},\qquad P(d)=d^2b(d).
\]
The Lean definition is \lean{Transpose.entryBits}, namely
{\small\texttt{Nat.log2 (d - 1) + 1}}; Lean's natural-number logarithm satisfies
{\small\texttt{Nat.log2 0 = 0}}.
A matrix is serialized in row-major order, with the bits of each entry
consecutive. The desired output uses the same bit order within each
transposed entry. These are \lean{Transpose.encode} and
\lean{Transpose.transpose}.

The statement $\WORD$ is \lean{Transpose.LinearTimeTransposable}: there
are one native machine $M$ and one constant $c\in\N$ such that every canonical
input of every dimension $d\geq1$ is transposed within $c(P(d)+1)$ transitions.
The statement $\BIT$ asks for one native machine and one constant that transpose
every square bit matrix of dyadic side $2^h$ in at most $c4^h$ transitions,
for every $h\geq0$, as in \lean{BitWord.DyadicBitOracle}.
These quantifiers require one machine for all sizes.

The compiler models in \cref{sec:geometry,sec:scans} are separate finite
models used to study possible reductions in the opposite direction. Their
operations are defined directly; there is no assertion that every native
computation has one of those forms. In particular, charging a word-oracle
call does not turn arbitrary surrounding computation into a permitted stage.
